% !TeX program = xelatex
% UTF-8 transcription of all four source pages, in photographic order.
% Compile here with: xelatex "cord disorders.tex"
% May also be \input into a document that already loads xepersian.
% Source wording (including clinical claims and English abbreviations) is retained.
% Bold text corresponds to yellow highlighting in the photographs.
% Uncertain handwriting is explicitly marked, rather than silently completed.
% pic01.jpg = 20261010_192203619.jpg
% pic02.jpg = 20261010_192209744.jpg
% pic03.jpg = 20261010_192213006.jpg
% pic04.jpg = 20261010_192217639.jpg
\ifdefined\settextfont
  \def\CordStandalone{0}
\else
  \def\CordStandalone{1}
  \documentclass[12pt,a4paper]{article}
  \usepackage[margin=2cm]{geometry}
  \usepackage{amsmath,amssymb,graphicx}
  \usepackage{xepersian}
  \settextfont{Amiri}
  \setlatintextfont{FreeSerif}
\fi
\newcommand{\corden}[1]{\lr{#1}}
\newcommand{\cordunclear}[1]{\textbf{[خوانش نامطمئن: #1]}}
\newcommand{\cordpage}[2]{\clearpage\section*{صفحهٔ #1}
  \noindent\corden{#2}\par
  \noindent\corden{Date: / /}\quad
  \corden{Sat. Sun. Mon. Tue. Thu. Wed. Fri.}\par
  \noindent\corden{Subject:} ۱۹ ـ بیماری‌های نخاع\par\medskip}
\ifnum\CordStandalone=1 \begin{document}\fi

\cordpage{۱}{pic01.jpg — 20261010\_192203619.jpg}
\section*{بیماری‌های نخاع}
\begin{itemize}
\item هایلایت بیماری‌های نخاع: وجود سطح حسی، حرکتی، اتونوم.
\item اختلال حرکتی در بیماری نخاع: افزایش رفلکس‌ها، بابینسکی، از بین رفتن رفلکس پوستی شکمی، اختلال اسفنکتر از نوع اسپاستیک.
\end{itemize}
\subsection*{ضایعات نخاع سرویکال}
\begin{itemize}
\item اگر خیلی نزدیک به \corden{pons} باشد: به‌علت مراکز کنترل تنفس، کشنده.
\item بخش فوقانی: کوادری‌پلژی اسپاستیک، ضعف دیافراگم، تنفس به کمک عضلات فرعی.
\item بخش تحتانی: در اندام فوقانی فلج شل / در اندام تحتانی فلج اسپاستیک.
\item سندرم هورنر (میوز، پتوز، انهیدروز)، اختلالات اسفنکتر.
\end{itemize}
\subsection*{ضایعات نخاع توراسیک}
\begin{itemize}
\item ضعف اندام‌های تحتانی به‌صورت پاراپارزی یا پاراپلژی.
\item اختلال حس و عملکرد مثانه و روده.
\item درماتوم نیپل (\corden{$T_4$})؛ درماتوم ناف (\corden{$T_{10}$}).
\end{itemize}
\subsection*{ضایعات نخاع لومبار}
\begin{itemize}
\item \corden{$L_2$--$L_4$}: فلج فلکشن و اداکشن ران؛ ضعف اکستنشن زانو؛ از بین رفتن رفلکس زانو.
\item \corden{$L_5$--$S_1$}: فلج حرکات پا؛ فلج \corden{Ankle flexion} زانو؛ فلج اکستنشن ران؛ از بین رفتن رفلکس آشیل.
\end{itemize}
\subsection*{ضایعات نخاع ساکرال و کونوس مدولاریس (سندرم کونوس)}
\begin{itemize}
\item احتباس ادرار، بی‌اختیاری همراه با شل شدن تون اسفنکتر مقعد.
\item اختلال حس اطراف مقعد (\corden{Saddle-anesth}).
\item ناتوانی جنسی.
\end{itemize}

\cordpage{۲}{pic02.jpg — 20261010\_192209744.jpg}
\begin{itemize}
\item در اختلالات این ناحیه قدرت عضلانی حفظ شده است.
\item ضایعات \corden{Cauda Equina}: درد قسمت تحتانی کمر، درد رادیکولر، ضعف، اختلال حس غیرقرینه در اندام تحتانی.
\end{itemize}
\subsection*{سندرم \corden{Hemicord} یا \corden{Brown-Sequard}}
\begin{itemize}
\item اختلال عملکرد و یا قطع یک نیمهٔ نخاع.
\item ضعف اندام‌ها در سمت ضایعه: درگیری راه‌های کورتیکواسپاینال.
\item از بین رفتن حس پوزیشن، ارتعاش و حس عمقی در سمت ضایعه: درگیری ستون خلفی.
\item از بین رفتن حس درد و حرارت در سمت مقابل ضایعه: درگیری اسپاینوتالامیک.
\item بابینسکی در سمت ضایعه.
\end{itemize}
\subsection*{سندرم \corden{Central Cord}}
\begin{itemize}
\item آسیب به سلول‌های عصبی مادهٔ خاکستری و مسیرهای متقاطع اسپاینوتالامیک در نزدیک کانال مرکزی نخاع.
\item علت: تروما ناشی از هایپراکستنشن گردن، سیرنگومیلی، تومورها، ایسکمی شریان قدامی نخاع. [یادداشت متصل به سیرنگومیلی: مهم‌ترین علت.]
\item علائم: ضعف در اندام فوقانی شدیدتر از اندام تحتانی است.
\item اختلال حس \corden{Dissociated}: از بین رفتن حس درد و حرارت و حفظ حس پوزیشن و ارتعاش.
\item علائم دوطرفه است و هر دو اندام فوقانی را درگیر می‌کند.
\end{itemize}
\subsection*{انفارکتوس نخاع}
\begin{itemize}
\item \textbf{۳ شریان} به نخاع خون‌رسانی می‌کنند: \textbf{۱ قدامی، ۲ خلفی}.
\end{itemize}

\cordpage{۳}{pic03.jpg — 20261010\_192213006.jpg}
\begin{itemize}
\item \textbf{دو سوم قدامی نخاع}: شاخه‌های \corden{penetrating} شریان قدامی؛ از \textbf{شریان ورتبرال}.
\item یک سوم خلفی نخاع: دو شریان خلفی.
\item \textbf{نخاع فوقانی توراسیک}: شایع‌ترین محل انفارکت به‌دلیل خون‌گیری ضعیف.
\item در بیمار با هایپوتنشن سیستمیک: \textbf{انفارکت \corden{$T_3$--$T_4$}} و مناطق حدفاصل خون‌رسانی شریان‌ها.
\item اتیولوژی: آترواسکلروز آئورت، آنوریسم دایسکان آئورت، کاهش \corden{BP}.
\item سایر علل: آمبولی قلبی، واسکولیت (\corden{SLE, APS})، پارگی آنوریسم آئورت.
\item علائم شریان قدامی: \textbf{پاراپلژی، کوادری‌پلژی، اختلال حس درد و حرارت و اختلال اسفنکتری}.
\item \textbf{حس ارتعاش و \corden{position} سالم است}.
\item درد شدید خط وسط و انتشار یابنده به پشت.
\item در آغاز به‌علت شوک نخاعی همهٔ رفلکس‌ها از بین می‌روند ولی با گذشت زمان هایپررفلکسی و اسپاستیسیته رخ می‌دهد.
\item علائم شریان‌های خلفی: \textbf{حس \corden{position} و ارتعاش از بین می‌رود}.
\item \textbf{تشخیص: \corden{MRI}}.
\item در موارد ناشی از ترومبوآمبولی آنتی‌کوآگولان اندیکاسیون ندارد؛ به‌غیر از مواردی که انفارکتوس کامل نگردیده و در حال پیشرفت است.
\item در \corden{APS}: آنتی‌کوآگولان ضروری است.
\end{itemize}
\subsection*{میلوپاتی تروماتیک}
\begin{itemize}
\item اتیولوژی: تروماهای نافذ، هایپرفلکشن یا هایپراکستنشن گردن، شکستگی در رفتگی ستون فقرات.
\item شدیدترین آسیب معمولاً به بخش مرکزی نخاع و به‌شکل سندرم سنترال کورد است.
\item علائم بالینی: ۱ ـ شوک نخاعی: از بین رفتن کامل حس، فلج کامل و سریع حرکتی، آرفلکسی در زیر سگمان آسیب‌دیده، فلج سیستم اتونوم.
\end{itemize}

\cordpage{۴}{pic04.jpg — 20261010\_192217639.jpg}
\begin{itemize}
\item فلج سیستم اتونوم: احتباس ادراری، ناتوانی جنسی، بی‌اختیاری مدفوع، از بین رفتن تون وازوموتور، تورم و سردی انتهاها، خشکی پوست.
\item ۲ ـ درگیری نورون محرکهٔ فوقانی: بعد از ۱ تا ۶ هفته.
\item ضعف، اسپاستیسیته، افزایش رفلکس‌ها، مثانهٔ اسپاستیک، بهبود نسبی انواع حس.
\end{itemize}
\subsection*{میلوپاتی ناشی از کمبود \corden{$B_{12}$} (دژنراسیون مرکب تحت حاد)}
\begin{latin}
\noindent Subacute Combined Degeneration
\end{latin}
\begin{itemize}
\item \corden{SCD}: ستون‌های خلفی و جانبی نخاع.
\item رژیم گیاه‌خواری مطلق.
\item علائم: \textbf{پارستزی دست‌ها و پاها}، از بین رفتن حس ارتعاش و پوزیشن.
\item \textbf{اختلال \corden{Gait}} (اسپاستیک و آتاکسیک).
\item در مراحل شدید: آتروفی عصب اپتیک، تحریک‌پذیری، اختلال خلق، دمانس، تغییرات رفتاری.
\item نشانهٔ رومبرگ: به‌دلیل درگیری ستون خلفی (آتاکسی حسی).
\item بابینسکی.
\item هایپو یا هایپررفلکسی تاندونی.
\item علامت لرمیت.
\item \cordunclear{گلوس؛ آتروفی؛ ماکروسیتیک — واژه‌ها و جداسازی آن‌ها در این سطر روشن نیست}.
\item تشخیص: سطح \corden{$B_{12}$} سرم / تغییر سیگنال ستون خلفی و طرفی در \corden{MRI}.
\item اگر زود درمان شود تمام علائم رفع می‌شوند.
\end{itemize}

% Preserve every original page for comparison with the transcription.
% Look both beside this file and relative to the repository root.
\begingroup
\graphicspath{{images/}{neurology/cord disorders/images/}}
\clearpage
\newcommand{\cordscan}[2]{\clearpage\subsection*{تصویر اصلی #1}
  \begin{center}
    \includegraphics[width=\linewidth,height=0.88\textheight,keepaspectratio]{#2}
  \end{center}}
\cordscan{۱}{pic01.jpg}
\cordscan{۲}{pic02.jpg}
\cordscan{۳}{pic03.jpg}
\cordscan{۴}{pic04.jpg}
\endgroup
\ifnum\CordStandalone=1
  \let\cordfinish\enddocument
\else
  \let\cordfinish\relax
\fi
\cordfinish
